\documentclass[a4paper,12pt]{article}
%\usepackage[lite,subscriptcorrection,slantedGreek,nofontinfo]{mtpro2}
%\usepackage{times}
\usepackage{amsfonts}
\usepackage{mathrsfs}
\usepackage{amsmath}
\usepackage{amssymb}
\usepackage{framed}
\usepackage[medium]{titlesec}
\usepackage{bm}
\usepackage{cite}
\usepackage[normalem]{ulem}
\usepackage{extarrows}
\usepackage{slashed}
\usepackage{isodateo}
\usepackage{graphicx}
\usepackage{xcolor}
\usepackage[bookmarksnumbered=true,bookmarksopen=true]{hyperref}
\hypersetup{colorlinks,%
            linkcolor=[rgb]{0,0.3,0.6},%
            citecolor=[rgb]{0,0.3,0.6},%
            urlcolor=[rgb]{0,0.3,0.6}}
\usepackage[hmargin=.7in,vmargin=1.1in]{geometry}
\usepackage{indentfirst}
\usepackage{booktabs}
\usepackage{tikz}
\usetikzlibrary{decorations.markings}

\linespread{1.1}
\newcommand{\FR}[2]{\displaystyle\frac{\,{#1}\,}{#2}}
\newcommand{\fr}[2]{\mbox{$\frac{\,{#1}\,}{#2}$}}
\newcommand{\n}{\nonumber}
\renewcommand{\rm}{\mathrm}
\renewcommand{\thefootnote}{\arabic{footnote}}

\newcommand{\red}{\color{red}}
\newcommand{\blue}{\color{blue}}
\newcommand{\gray}{\color{gray}}
\newcommand{\hl}{\color[rgb]{1,0,.5}}

\def\bge{\begin{equation}}
\def\ede{\end{equation}}
\def\bga{\begin{aligned}}
\def\eda{\end{aligned}}
\def\bgb{\begin{bmatrix}}
\def\edb{\end{bmatrix}}
\def\bgp{\begin{pmatrix}}
\def\edp{\end{pmatrix}}
\def\bgs{\begin{subequations}}
\def\eds{\end{subequations}}
\newcommand{\order}[1]{\mathcal{O}({#1})}
\def\di{{\mathrm{d}}}
\def\mb{\mathbf}
\def\ms{\mathscr}
\def\mc{\mathcal}
\def\pd{\partial}
\def\la{\langle}\def\ra{\rangle}
\def\tr{\mathrm{\,tr\,}}
\def\ii{\mathrm{i}}
\def\de{\delta}
\def\ep{\epsilon}
\def\lam{\lambda}
\def\si{\sigma}
\def\ka{\kappa}
\def\tn{\widetilde\nu}
\def\be{\beta}
\def\Re{\mathrm{Re}\,}
\def\Im{\mathrm{Im}\,}

\usepackage{mdframed}
\newmdenv[backgroundcolor=gray!15,%
          skipabove=0pt,%
          skipbelow=5pt,%
          leftmargin=0pt,%
          rightmargin=0pt,%
          innertopmargin=-5pt,%
          innerbottommargin=7pt,%
          innerleftmargin=2pt,%
          innerrightmargin=2pt,%
          splittopskip=0pt,%
          splitbottomskip=0pt,%
          linewidth=0pt,%
          nobreak=true]{keyeqn}

\titleformat{\section}
{\normalfont\fontsize{15}{20}\bfseries}{\thesection}{1em}{}

\usepackage[T1]{fontenc}
\usepackage{lmodern}

\newcommand{\ob}[1]{\mkern 2mu \overline{\mkern -2mu #1 \mkern -2mu}\mkern 2mu}
\newcommand{\wt}[1]{\mkern 2mu \widetilde{\mkern -2mu #1 \mkern -2mu}\mkern 2mu}
\newcommand{\wh}[1]{\mkern 2mu \widehat{\mkern-2mu#1\mkern-2mu}\mkern 2mu}
\tikzset{
 propforward/.style={postaction={decorate},decoration={markings,
 mark=at position .53 with {\arrow{>}}}},
 propreverse/.style={postaction={decorate},decoration={markings,
 mark=at position .47 with {\arrow{<}}}},
 propinward/.style={postaction={decorate},decoration={markings,
 mark=at position .31 with {\arrow{>}},
 mark=at position .69 with {\arrow{<}}}},
 propoutward/.style={postaction={decorate},decoration={markings,
 mark=at position .31 with {\arrow{<}},
 mark=at position .69 with {\arrow{>}}}}
}
\newcommand{\propedge}[3]{%
 \ifcase#1\relax
 \or\draw[propforward] (#2)--(#3);
 \or\draw[propreverse] (#2)--(#3);
 \or\draw[propinward] (#2)--(#3);
 \or\draw[propoutward] (#2)--(#3);
 \fi}
\newcommand{\weylbox}[9]{%
\begin{tikzpicture}[baseline=(current bounding box.center),scale=.62]
 \coordinate (v1) at (-.62,.62);
 \coordinate (v2) at (.62,.62);
 \coordinate (v3) at (.62,-.62);
 \coordinate (v4) at (-.62,-.62);
 \propedge{#6}{v1}{v2}
 \begin{scope}[blue!65!black,very thick]
  \propedge{#7}{v2}{v3}
  \propedge{#9}{v4}{v1}
 \end{scope}
 \propedge{#8}{v3}{v4}
 \draw (v1)--(-.95,.95) (v2)--(.95,.95)
       (v3)--(.95,-.95) (v4)--(-.95,-.95);
 \fill (v1) circle (1.4pt) (v2) circle (1.4pt)
       (v3) circle (1.4pt) (v4) circle (1.4pt);
 \node[above left,font=\scriptsize] at (v1) {$#2$};
 \node[above right,font=\scriptsize] at (v2) {$#3$};
 \node[below right,font=\scriptsize] at (v3) {$#4$};
 \node[below left,font=\scriptsize] at (v4) {$#5$};
 \node[font=\scriptsize] at (0,-1.22) {#1};
\end{tikzpicture}}

\title{\Large\textbf{Top Quark Signal in the Cosmological Higgs Collider}\\[2mm]}
\author{}

\begin{document}

\date{}
\maketitle

\section{Setup}

We calculate the top-quark contribution to the connected Higgs four-point function in the
collapsed, or $s$-channel squeezed, limit. 
As in the gauge-boson calculation, the Higgs is treated as massless during inflation.  
Its Schwinger--Keldysh bulk-to-boundary propagator is:
\begin{align}
    G_a(k;\tau)=\frac{H^2}{2k^3}(1-\ii a k\tau)e^{\ii a k\tau},
\end{align}
and the scale factor is $a(\tau)=-(H\tau)^{-1}$.

In unitary gauge the top Yukawa interaction is
\begin{align}
    \mathscr L_t
    &=\ii t^\dagger\bar\si^\mu\pd_\mu t
      +\ii t^{c\dagger}\bar\si^\mu\pd_\mu t^c
      -a(\tau)\big[m_t+g_t\de h\big]
      \big(tt^c+t^\dagger t^{c\dagger}\big),
      \label{eq:dirac-L}\\
    g_t&\equiv\frac{y_t}{\sqrt2},\qquad
    m_t=g_t h_0,\qquad
    \tn\equiv\frac{m_t}{H}.
\end{align}
Here $t$ and $t^c$ are two left-handed Weyl fields and have already been conformally rescaled by $a^{3/2}$.  
Hence every Yukawa vertex carries one power of $a(\tau)$, rather than $a^4(\tau)$.  

\subsection{From the Dirac top to two Majorana spinors}

The fermion formalism of Ref.~\cite{you_rhn} is written for a two-component Majorana field.  
It can be applied to the top without changing the physics by the unitary Takagi rotation
\begin{align}
    t=\frac{N_1+\ii N_2}{\sqrt2},\qquad
    t^c=\frac{N_1-\ii N_2}{\sqrt2}.
    \label{eq:takagi}
\end{align}
Because the contraction of two Weyl spinors is symmetric,
$N_1N_2=N_2N_1$, this gives
\begin{align}
    tt^c=\frac12(N_1N_1+N_2N_2).
\end{align}
Thus Eq.~\eqref{eq:dirac-L} becomes
\begin{align}
    \mathscr L_t=\sum_{A=1}^{2}\left[
    \ii N_A^\dagger\bar\si^\mu\pd_\mu N_A
    -\frac{a m_t}{2}(N_AN_A+N_A^\dagger N_A^\dagger)
    -\frac{a g_t}{2}\de h(N_AN_A+N_A^\dagger N_A^\dagger)
    \right].
    \label{eq:maj-L}
\end{align}
Both $N_A$ therefore have the same positive Majorana mass and the same Higgs coupling.
This is preferable to diagonalizing the mass matrix with eigenvalues $\pm m_t$: the phase
in Eq.~\eqref{eq:takagi} puts the result directly in the positive-mass convention used in
Ref.~\cite{you_rhn}.

There is no additional factor of two in a closed top loop.  A single Majorana functional
determinant is one half of a Dirac determinant, while the sum over $A=1,2$ supplies the
opposite factor two.  The net multiplicity is consequently the standard color factor
$N_c=3$, together with the usual minus sign for a closed fermion loop.

For momentum $\mb k=k(\sin\theta\cos\phi,\sin\theta\sin\phi,\cos\theta)$, we use the
helicity eigenspinors
\begin{align}
 h_+(\theta,\phi)&=
 \begin{pmatrix}e^{-\ii\phi/2}\cos(\theta/2)\\
 e^{\ii\phi/2}\sin(\theta/2)\end{pmatrix},&
 h_-(\theta,\phi)&=
 \begin{pmatrix}-e^{-\ii\phi/2}\sin(\theta/2)\\
 e^{\ii\phi/2}\cos(\theta/2)\end{pmatrix},
 \label{eq:helicity-spinors}
\end{align}
The mode functions:
\begin{align}
 u_+(k,\tau)&=\frac{\tn}{\sqrt{-2k\tau}}
 W_{-1/2,\ii\tn}(2\ii k\tau),&
 u_-(k,\tau)&=\frac{1}{\sqrt{-2k\tau}}
 W_{1/2,\ii\tn}(2\ii k\tau),
 \\
 v_+(k,\tau)&=\frac{1}{\sqrt{-2k\tau}}
 W_{1/2,\ii\tn}(2\ii k\tau),&
 v_-(k,\tau)&=\frac{\tn}{\sqrt{-2k\tau}}
 W_{-1/2,\ii\tn}(2\ii k\tau).
 \label{eq:modes}
\end{align}

\subsection{Fermion propagators and their nonlocal parts}

There are three useful two-component propagators,
\begin{align}
 D_{ab\,\alpha\dot\be}&\equiv
 \la N_\alpha N^\dagger_{\dot\be}\ra_{ab},&
 \widehat D_{ab\,\alpha\be}&\equiv
 \la N_\alpha N_\be\ra_{ab},&
 \check D_{ab\,\dot\alpha\dot\be}&\equiv
 \la N^\dagger_{\dot\alpha}N^\dagger_{\dot\be}\ra_{ab}.
\end{align}
Following Ref.~\cite{you_rhn}, we call $D$, $\widehat D$, and $\check D$ the type-1,
type-2, and type-3 propagators, respectively.  
The nonlocal late-time parts of the symmetric type-2 and type-3
propagators are especially simple:
\begin{align}
 \left[\widehat D^{(S)\,\alpha}{}_{\be}
 (k;\tau_1,\tau_2)\right]_{\rm NLoc}
 &=\sum_{f=\pm}\ka_f(\tn)
 (4k^2\tau_1\tau_2)^{\ii f\tn}\de^\alpha_{\ \be},
 \\
 \left[\check D^{(S)\,\dot\alpha}{}_{\dot\be}
 (k;\tau_1,\tau_2)\right]_{\rm NLoc}
 &=\sum_{f=\pm}\ka_f(\tn)
 (4k^2\tau_1\tau_2)^{\ii f\tn}\de^{\dot\alpha}_{\ \dot\be},
 \label{eq:soft-type23}
\end{align}
where
\begin{align}
 \ka_f(\tn)
 =\frac{\tn\Gamma^2(-2\ii f\tn)}
 {\Gamma(1-\ii f\tn)\Gamma(-\ii f\tn)},
 \qquad \ka_- =\ka_+^*.
 \label{eq:kappa}
\end{align}
At zero chemical potential the helicity tensor in the unraised propagator reduces to the
invariant antisymmetric tensor,
\begin{align}
 h_{+\alpha}h_{-\be}-h_{-\alpha}h_{+\be}=\ep_{\alpha\be}.
\end{align}
It follows that the two Wightman orderings have the same type-2 nonlocal part. 

The type-1 propagator behaves differently.  At zero chemical potential its symmetric
nonlocal part vanishes, while its antisymmetric part is proportional to
$\mb\si\cdot\wh{\mb k}$.  Consequently a pair of type-1 soft lines is not covered by the
factorized symmetric prescription of Ref.~\cite{you_rhn}.  We return to this point in
Sec.~\ref{sec:antisym}.

\section{Box diagram and classification of the sixteen contractions}
\label{sec:sixteen}

For one cyclic ordering of the external Higgs legs, the box has the schematic SK form
\begin{align}
 \la\de h^4\ra'_{\rm Box}
 =-N_c(\ii g_t)^4\sum_{a,b,c,d=\pm}(abcd)
 \int\prod_{j=1}^{4}\di\tau_j\,a(\tau_j)G_{a_j}(k_j;\tau_j)
 \int\frac{\di^3\mb q}{(2\pi)^3}\,\mathcal{S}_{abcd},
 \label{eq:box-start}
\end{align}
where $(a_1,a_2,a_3,a_4)=(a,b,c,d)$ and $\mathcal S_{abcd}$ denotes the cyclic spinor
contraction of the four fermion propagators.  Crossed cyclic orderings are obtained by
permuting the external labels.

The nonlocal $s$-channel signal arises when the lines with momenta $\mb q$ and
$\mb K=\mb q+\mb k_s$ are soft.
At every Yukawa vertex one may choose either $NN$ or $N^\dagger N^\dagger$.  
We denote these choices by $N$ and $\bar N$.  
They give $2^4=16$ Weyl contractions for a fixed cyclic ordering of the four external Higgs legs.  
These are terms in the Majorana decomposition of one box topology, not sixteen different momentum topologies.
For compactness, on an oriented edge we write
\begin{align}
 NN:\widehat D,\qquad
 N\bar N:D,\qquad
 \bar NN:\bar D,\qquad
 \bar N\bar N:\check D .
 \label{eq:edge-rule}
\end{align}
The reversed type-1 propagator is not an additional independent object:
\begin{align}
 \bar D_{ab}^{\dot\alpha\be}(\mb k;\tau,\tau')
 \equiv-D_{ba}^{\be\dot\alpha}(-\mb k;\tau',\tau).
 \label{eq:Dbar}
\end{align}
This relation fixes the Fermi sign and the SK/time reversal for every $\bar D$ entry.
The hard edges are $(12)$ and $(34)$, and the soft edges are $(23)$ and $(41)$.  Before
making either a symmetric or an antisymmetric soft-line replacement, the complete list is
\begin{center}
\small
\begin{tabular}{ccllc}
\toprule
No. & $(r_1r_2r_3r_4)$ & $(12,23,34,41)$ & soft class & $n_1$\\
\midrule
1  & $NNNN$                         & $(\widehat D,\widehat D,\widehat D,\widehat D)$ & $AA$ &0\\
2  & $NNN\bar N$                    & $(\widehat D,\widehat D,D,\bar D)$ & $AN$ &1\\
3  & $NN\bar NN$                    & $(\widehat D,D,\bar D,\widehat D)$ & $NA$ &1\\
4  & $NN\bar N\bar N$               & $(\widehat D,D,\check D,\bar D)$ & $NN$ &2\\
5  & $N\bar NNN$                    & $(D,\bar D,\widehat D,\widehat D)$ & $NA$ &1\\
6  & $N\bar NN\bar N$               & $(D,\bar D,D,\bar D)$ & $NN$ &2\\
7  & $N\bar N\bar NN$               & $(D,\check D,\bar D,\widehat D)$ & $AA$ &0\\
8  & $N\bar N\bar N\bar N$          & $(D,\check D,\check D,\bar D)$ & $AN$ &1\\
9  & $\bar NNNN$                    & $(\bar D,\widehat D,\widehat D,D)$ & $AN$ &1\\
10 & $\bar NNN\bar N$               & $(\bar D,\widehat D,D,\check D)$ & $AA$ &0\\
11 & $\bar NN\bar NN$               & $(\bar D,D,\bar D,D)$ & $NN$ &2\\
12 & $\bar NN\bar N\bar N$          & $(\bar D,D,\check D,\check D)$ & $NA$ &1\\
13 & $\bar N\bar NNN$               & $(\check D,\bar D,\widehat D,D)$ & $NN$ &2\\
14 & $\bar N\bar NN\bar N$          & $(\check D,\bar D,D,\check D)$ & $NA$ &1\\
15 & $\bar N\bar N\bar NN$          & $(\check D,\check D,\bar D,D)$ & $AN$ &1\\
16 & $\bar N\bar N\bar N\bar N$     & $(\check D,\check D,\check D,\check D)$ & $AA$ &0\\
\bottomrule
\end{tabular}
\end{center}
\begin{figure}[ht]
\centering
\begin{tabular}{cccc}
\weylbox{1:\ $AA$}{N}{N}{N}{N}{3}{3}{3}{3}&
\weylbox{2:\ $AN$}{N}{N}{N}{\bar N}{3}{3}{1}{2}&
\weylbox{3:\ $NA$}{N}{N}{\bar N}{N}{3}{1}{2}{3}&
\weylbox{4:\ $NN$}{N}{N}{\bar N}{\bar N}{3}{1}{4}{2}\\[3mm]
\weylbox{5:\ $NA$}{N}{\bar N}{N}{N}{1}{2}{3}{3}&
\weylbox{6:\ $NN$}{N}{\bar N}{N}{\bar N}{1}{2}{1}{2}&
\weylbox{7:\ $AA$}{N}{\bar N}{\bar N}{N}{1}{4}{2}{3}&
\weylbox{8:\ $AN$}{N}{\bar N}{\bar N}{\bar N}{1}{4}{4}{2}\\[3mm]
\weylbox{9:\ $AN$}{\bar N}{N}{N}{N}{2}{3}{3}{1}&
\weylbox{10:\ $AA$}{\bar N}{N}{N}{\bar N}{2}{3}{1}{4}&
\weylbox{11:\ $NN$}{\bar N}{N}{\bar N}{N}{2}{1}{2}{1}&
\weylbox{12:\ $NA$}{\bar N}{N}{\bar N}{\bar N}{2}{1}{4}{4}\\[3mm]
\weylbox{13:\ $NN$}{\bar N}{\bar N}{N}{N}{4}{2}{3}{1}&
\weylbox{14:\ $NA$}{\bar N}{\bar N}{N}{\bar N}{4}{2}{1}{4}&
\weylbox{15:\ $AN$}{\bar N}{\bar N}{\bar N}{N}{4}{4}{2}{1}&
\weylbox{16:\ $AA$}{\bar N}{\bar N}{\bar N}{\bar N}{4}{4}{4}{4}
\end{tabular}
\caption{The sixteen Weyl contractions for one cyclic box, using the fermion-arrow
convention of Ref.~\cite{you_rhn}.  A type-1 line has one arrow directed from $N$ to
$N^\dagger$; the two arrows on a type-2 line point inward, while those on a type-3 line
point outward.  The horizontal edges $(12)$ and $(34)$ are hard; the blue vertical edges
$(23)$ and $(41)$ are soft.  The label below each diagram is its soft class.}
\label{fig:sixteen-boxes}
\end{figure}
Here $A$ is only a compact class label for a type-2 $\widehat D$ or type-3 $\check D$
soft line, while $N$ denotes a type-1 $D$ or $\bar D$ soft line; $n_1$ counts the
type-1 soft lines.  

\section{Symmetric soft-line contribution}

We first replace both soft lines by their symmetric nonlocal parts.  At zero chemical
potential,
\begin{align}
 [D^{(S)}]_{\rm NLoc}=0,\qquad
 [\widehat D^{(S)}]_{\rm NLoc}\ne0,\qquad
 [\check D^{(S)}]_{\rm NLoc}\ne0.
 \label{eq:symmetric-zero}
\end{align}
Consequently
\begin{align}
 C^{(S)}_\ell&=0,\qquad
 \ell\in\{2,3,4,5,6,8,9,11,12,13,14,15\},\\
 {\cal G}_{AA}&=\{1,7,10,16\}.
 \label{eq:GAA}
\end{align}
Thus twelve rows vanish and only four rows survive.  This condition concerns the
\emph{soft} edges only: the hard edges of rows 1, 7, 10, and 16 are respectively
$(2,2)$, $(1,1)$, $(1,1)$, and $(3,3)$.

\subsection{Loop momentum integral and bubble subgraph}

Following the gauge-boson calculation, we keep the two hard propagators unevaluated and
perform the soft loop momentum first.  For a surviving word
$w_\ell=(r_1r_2r_3r_4)$, define
\begin{align}
 {\cal H}_{\ell,abcd}(\tau_i)
 =\tr_2\!\left[
 {\cal D}^{r_1r_2}_{ab}(\mb k_1;\tau_1,\tau_2)
 \Delta^{r_2r_3}
 {\cal D}^{r_3r_4}_{cd}(\mb k_3;\tau_3,\tau_4)
 \Delta^{r_4r_1}\right],
 \label{eq:hard-kernel-before-time}
\end{align}
where the raised-index type-2/type-3 soft matrices are
$\Delta^{NN}=\Delta^{\bar N\bar N}=\si^0$.  Before doing the $\mb q$ integral,
\begin{align}
 {\cal LOOP}^{(S)}_{\ell}
 &=\sum_{f=\pm}\ka_f^2\,2^{4\ii f\tn}
 \left[\prod_{j=1}^{4}(-\tau_j)^{\ii f\tn}\right]
 {\cal H}_{\ell,abcd}(\tau_i)
 \int\frac{\di^3\mb q}{(2\pi)^3}q^{2\ii f\tn}K^{2\ii f\tn}.
 \label{eq:symmetric-loop-before-q}
\end{align}
The loop integral therefore separates into a bubble subgraph and one extra time factor
at each vertex:
\begin{align}
 {\cal LOOP}^{(S)}_{\ell}
 =\sum_{f=\pm}
 \left[\prod_{j=1}^{4}(-\tau_j)^{\ii f\tn}\right]
 {\cal B}^{(S)}_f(k_s)\,{\cal H}_{\ell,abcd}(\tau_i),
 \label{eq:symmetric-loop-after-q}
\end{align}
where
\begin{align}
 \mathcal B_f^{(S)}(k_s)
 &=\ka_f^2 2^{4\ii f\tn}
 \int\frac{\di^3\mb q}{(2\pi)^3}q^{2\ii f\tn}K^{2\ii f\tn}\\
 &=\ka_f^2 2^{4\ii f\tn}
 \frac{k_s^{3+4\ii f\tn}}{(4\pi)^{3/2}}
 f_0(-\ii f\tn,-\ii f\tn).
 \label{eq:soft-bubble}
\end{align}
In the notation of the gauge-boson calculation,
\begin{align}
 f_0(x,y)=
 \Gamma\left[\begin{array}{c}
 x+y-\frac32,\ \frac32-x,\ \frac32-y\\
 3-x-y,\ x,\ y
 \end{array}\right].
 \label{eq:f0}
\end{align}
The factor $2^{4\ii f\tn}$ follows directly from the normalization in
Eq.~\eqref{eq:soft-type23}.

\subsection{Left and right subgraphs}

The extra powers in Eq.~\eqref{eq:symmetric-loop-after-q} are now absorbed into the
left and right time integrals.  We first write these subgraphs before introducing any
seed.  For the chirality word $w_\ell=(r_1r_2r_3r_4)$, their matrix-valued integral
definitions are
\begin{subequations}\label{eq:T-time-def}
\begin{align}
 \mathcal T_{f,\ell}^{(L)}
 &=(\ii g_t)^2\sum_{a,b=\pm}ab
 \int_{-\infty}^{0}\di\tau_1\di\tau_2\,
 a(\tau_1)a(\tau_2)G_a(k_1;\tau_1)G_b(k_2;\tau_2)\nonumber\\
 &\hspace{14mm}\times
 (-\tau_1)^{\ii f\tn}(-\tau_2)^{\ii f\tn}
 {\cal D}^{r_1r_2}_{ab}(\mb k_1;\tau_1,\tau_2),\\
 \mathcal T_{f,\ell}^{(R)}
 &=(\ii g_t)^2\sum_{c,d=\pm}cd
 \int_{-\infty}^{0}\di\tau_3\di\tau_4\,
 a(\tau_3)a(\tau_4)G_c(k_3;\tau_3)G_d(k_4;\tau_4)\nonumber\\
 &\hspace{14mm}\times
 (-\tau_3)^{\ii f\tn}(-\tau_4)^{\ii f\tn}
 {\cal D}^{r_3r_4}_{cd}(\mb k_3;\tau_3,\tau_4).
\end{align}
\end{subequations}
The SK weights $ab$ and $cd$ are included here, whereas the soft coefficient
$\ka_f^2$ belongs to the bubble subgraph.  Inserting the scale factor and the two
massless Higgs wavefunctions into the first line gives
\begin{align}
 \mathcal T_{f,\ell}^{(L)}
 &=(\ii g_t)^2\frac{H^2}{4k_1^3k_2^3}
 \sum_{a,b=\pm}ab\int_{-\infty}^{0}\di\tau_1\di\tau_2\,
 (-\tau_1)^{p_f}(-\tau_2)^{p_f}\nonumber\\
 &\quad\times(1-\ii ak_1\tau_1)(1-\ii bk_2\tau_2)
 e^{\ii ak_1\tau_1+\ii bk_2\tau_2}
 {\cal D}^{r_1r_2}_{ab}(\mb k_1;\tau_1,\tau_2),
 \qquad p_f=-1+\ii f\tn .
 \label{eq:TL-before-seed}
\end{align}
The right subgraph is obtained by
$(1,2,a,b,r_1,r_2)\to(3,4,c,d,r_3,r_4)$.

Now define the two hard scales
\begin{align}
 k_L\equiv\frac{k_{12}}2,\qquad k_R\equiv\frac{k_{34}}2.
\end{align}
At leading order in the collapsed limit,
$k_1=k_2=k_L+\order(k_s)$ and $k_3=k_4=k_R+\order(k_s)$.
With $z_1=-k_L\tau_1$ and $z_2=-k_L\tau_2$,
Eq.~\eqref{eq:TL-before-seed} becomes
\begin{align}
 \mathcal T_{f,\ell}^{(L)}
 &=(\ii g_t)^2\frac{H^2}{4k_1^3k_2^3}
 k_L^{-2\ii f\tn}\sum_{a,b=\pm}ab
 \int_0^\infty\di z_1\di z_2\,
 z_1^{p_f}z_2^{p_f}(1+\ii az_1)(1+\ii bz_2)\nonumber\\
 &\hspace{28mm}\times e^{-\ii az_1-\ii bz_2}
 {\cal D}^{r_1r_2}_{ab}(\wh{\mb k}_1;z_1,z_2)
 +\order(k_s/k_i).
 \label{eq:TL-dimensionless}
\end{align}
The dimensionless propagator on the last line is the one used in
Ref.~\cite{you_rhn}; conformal rescaling makes its momentum dependence a function of
$k_L\tau_{1,2}=-z_{1,2}$ and the direction $\wh{\mb k}_1$.  Retaining unequal
$k_1/k_L$ and $k_2/k_L$ would merely give energy-deformed seeds and contributes beyond
the leading collapsed result considered here.

We can now read off the You seeds.  For endpoint chiralities
$r,s\in\{N,\bar N\}$ define
\begin{align}
 \mathcal I_{ab}^{rs;p_1p_2}(\wh{\mb k})
 =\int_0^\infty\di z_1\di z_2\,
 e^{-\ii a z_1-\ii b z_2}z_1^{p_1}z_2^{p_2}
 {\cal D}_{ab}^{rs}(\wh{\mb k};z_1,z_2).
 \label{eq:seed-def}
\end{align}
Here $\mathcal D^{NN}=\widehat D$, $\mathcal D^{N\bar N}=D$,
$\mathcal D^{\bar NN}=\bar D$, and $\mathcal D^{\bar N\bar N}=\check D$.
Local end-point terms are defined by analytic continuation and do not affect the
nonanalytic signal.  Expanding the two Higgs polynomials in
Eq.~\eqref{eq:TL-dimensionless} gives exactly
\begin{align}
 \mathcal H^{rs}_{f;ab}
 &=\mathcal I_{ab}^{rs;p_fp_f}
 +\ii a\mathcal I_{ab}^{rs;p_f+1,p_f}
 +\ii b\mathcal I_{ab}^{rs;p_f,p_f+1}
 -ab\mathcal I_{ab}^{rs;p_f+1,p_f+1},\\
 \mathcal Q_f^{rs}(\wh{\mb k})
 &=\sum_{a,b=\pm}ab\,\mathcal H^{rs}_{f;ab}(\wh{\mb k}).
 \label{eq:Q-def}
\end{align}
Therefore the two original time integrals reduce to
\begin{align}
 \mathcal T_{f,\ell}^{(L)}
 &=(\ii g_t)^2\frac{H^2}{4k_1^3k_2^3}
 \left(\frac{k_{12}}{2}\right)^{-2\ii f\tn}
 \mathcal Q_f^{r_1r_2}(\wh{\mb k}_1),\\
 \mathcal T_{f,\ell}^{(R)}
 &=(\ii g_t)^2\frac{H^2}{4k_3^3k_4^3}
 \left(\frac{k_{34}}{2}\right)^{-2\ii f\tn}
 \mathcal Q_f^{r_3r_4}(\wh{\mb k}_3).
 \label{eq:hard-subgraphs}
\end{align}
Appendix~\ref{app:seeds} gives the type-1, type-2, and type-3 seed matrices required by
Eq.~\eqref{eq:Q-def}.  Thus no hard time integral remains.  For the type-1 hard edges of
rows 7 and 10 we abbreviate
\begin{align}
 \mathcal Q_f^{N\bar N}(\wh{\mb k})
 =Q_{f,0}\si^0+Q_{f,1}\wh{\mb k}\cdot\mb\si .
 \label{eq:Q-rotation}
\end{align}

\subsection{Factorized form before the spinor contraction}

Let $w_\ell=(r_1r_2r_3r_4)$ denote the chirality word in row $\ell$.  After the
bubble integral has been done, the symmetric contribution takes exactly the same
left--bubble--right form as the gauge-boson box,
\begin{align}
 \left.\la\de h^4\ra'_{\rm Box}\right|_{S}
 =-N_c\sum_{\ell\in{\cal G}_{AA}}\sum_{f=\pm}
 {\cal T}^{(L)}_{f,\ell}\,{\cal B}^{(S)}_f\,
 {\cal T}^{(R)}_{f,\ell}\big|_{\text{spinor contraction}} .
 \label{eq:symmetric-factorization}
\end{align}
The scalar bubble coefficient is Eq.~\eqref{eq:soft-bubble}; all spinor indices reside in
the two hard blocks.  Since a symmetric type-2 or type-3 soft propagator becomes a
Kronecker delta after raising the index contracted at the adjacent Yukawa vertex, no
direction associated with $\mb q$ or $\mb K$ remains in this sector.  The spinor
contraction can therefore be completed before doing either hard time integral.

\subsection{Four spinor contractions and the closed symmetric result}
\label{sec:symmetric-four}

We now evaluate rows $1,7,10,16$ separately.  This also fixes the relative sign between
the type-1 and type-2/type-3 hard lines.  Let
\begin{align}
 E_{\alpha\be}=\ep_{\alpha\be},\qquad
 \widehat{\mb Q}_{f}=\widehat{\mathcal Q}_{f}E^{-1},\qquad
 \check{\mb Q}_{f}=E^{-1}\check{\mathcal Q}_{f},
 \label{eq:raised-hard-blocks}
\end{align}
where the placement of $E^{-1}$ is fixed by whether the indices are undotted or dotted.
These are just the mixed-index matrices obtained by contracting the hard propagator with
the two Yukawa tensors.  In the $\mb k\parallel\wh{\mb z}$ frame define their four pairs
of radial coefficients by
\begin{align}
 \widehat{\mb Q}_{+}(\wh{\mb z})
 &=\widehat q_0\si^0+\widehat q_1\si^3,&
 \check{\mb Q}_{+}(\wh{\mb z})
 &=\check q_0\si^0+\check q_1\bar\si^3,\nonumber\\
 \mb Q_{+}(\wh{\mb z})
 &=q_0\si^0+q_1\si^3,&
 \ob{\mb Q}_{+}(\wh{\mb z})
 &=\bar q_0\si^0+\bar q_1\bar\si^3,
 \qquad \bar\si^i=-\si^i .
 \label{eq:q-radial-def}
\end{align}
Here $\mb Q$ and $\ob{\mb Q}$ are respectively the $D$ and $\bar D$ hard blocks.
The eight quantities in Eq.~\eqref{eq:q-radial-def} depend only on $\tn$ at the leading
collapsed order.  The rotation rule of Ref.~\cite{you_rhn} gives
\begin{align}
 \si^3\longrightarrow\wh{\mb k}\cdot\mb\si,
 \qquad
 \bar\si^3\longrightarrow\wh{\mb k}\cdot\bar{\mb\si}.
 \label{eq:q-rotation}
\end{align}
We choose the same angles as in the gauge-boson box,
\begin{align}
 \mb k_s&=k_s(0,0,1),\nonumber\\
 \mb k_1&=k_1(\sin\theta_1,0,\cos\theta_1),\nonumber\\
 \mb k_3&=k_3(\sin\theta_3\cos\phi_3,
 \sin\theta_3\sin\phi_3,\cos\theta_3),
\end{align}
and use the single invariant
\begin{align}
 c_{13}\equiv\wh{\mb k}_1\cdot\wh{\mb k}_3
 =\cos\theta_1\cos\theta_3
 +\sin\theta_1\sin\theta_3\cos\phi_3 .
 \label{eq:angles}
\end{align}

The two raised soft lines are identities.  Using
$\tr(\si^i\si^j)=2\de^{ij}$ and
$\tr(\si^i\bar\si^j)=-2\de^{ij}$, the four contractions are
\begin{subequations}\label{eq:four-symmetric-traces}
\begin{align}
 {\cal W}_{+,1}
 &=\tr_2\!\left[
 \widehat{\mb Q}_{+}(\wh{\mb k}_1)
 \widehat{\mb Q}_{+}(\wh{\mb k}_3)\right]
 =2\left(\widehat q_0^{\,2}+\widehat q_1^{\,2}c_{13}\right),
 \label{eq:trace-row1}\\
 {\cal W}_{+,7}
 &=\tr_2\!\left[
 \mb Q_{+}(\wh{\mb k}_1)\ob{\mb Q}_{+}(\wh{\mb k}_3)\right]
 =2\left(q_0\bar q_0-q_1\bar q_1c_{13}\right),
 \label{eq:trace-row7}\\
 {\cal W}_{+,10}
 &=\tr_2\!\left[
 \ob{\mb Q}_{+}(\wh{\mb k}_1)\mb Q_{+}(\wh{\mb k}_3)\right]
 =2\left(\bar q_0q_0-\bar q_1q_1c_{13}\right),
 \label{eq:trace-row10}\\
 {\cal W}_{+,16}
 &=\tr_2\!\left[
 \check{\mb Q}_{+}(\wh{\mb k}_1)
 \check{\mb Q}_{+}(\wh{\mb k}_3)\right]
 =2\left(\check q_0^{\,2}+\check q_1^{\,2}c_{13}\right).
 \label{eq:trace-row16}
\end{align}
\end{subequations}
Thus rows 7 and 10 are equal after the radial coefficients have been extracted, but both
must be retained: they are the two opposite fermion-flow assignments.  Equations
\eqref{eq:four-symmetric-traces} also show directly why the symmetric angular dependence
cannot contain $\cos2\phi_3$, $\cos\theta_1$, or $\cos\theta_3$ separately.

For completeness, we now write the hard quantities in terms of the integral seeds of
Ref.~\cite{you_rhn}.  With $p=-1+\ii\tn$, define for
$X=D,\bar D,\widehat D,\check D$
\begin{align}
 \mb Q^{X}_{+}(\wh{\mb k})
 &=\sum_{a,b=\pm}ab\Big[
 \mb I^{X;p,p}_{ab}
 +\ii a\mb I^{X;p+1,p}_{ab}
 +\ii b\mb I^{X;p,p+1}_{ab}
 -ab\mb I^{X;p+1,p+1}_{ab}\Big](\wh{\mb k}).
 \label{eq:Yukawa-dressed-seed}
\end{align}
This is the scalar-Yukawa analogue of the subdiagram construction in
Ref.~\cite{you_rhn}: the four terms $(m,n)=(0,0),(1,0),(0,1),(1,1)$ come from the two
massless Higgs polynomials.  For $X=\widehat D,\check D$, the index-raising operation in
Eq.~\eqref{eq:raised-hard-blocks} is understood on the right-hand side.  More explicitly,
at leading collapsed order we may set
$k_1=k_2=k_{12}/2+\order(k_s)$ and
$k_3=k_4=k_{34}/2+\order(k_s)$.  Putting
\begin{align}
 {\cal N}_L&=(\ii g_t)^2\frac{H^2}{4k_1^3k_2^3}
 \left(\frac{k_{12}}2\right)^{-2\ii\tn},&
 {\cal N}_R&=(\ii g_t)^2\frac{H^2}{4k_3^3k_4^3}
 \left(\frac{k_{34}}2\right)^{-2\ii\tn},
 \label{eq:NLNR}
\end{align}
the four pairs of left and right subdiagrams are
\begin{align}
 (\mathcal T^{(L)}_{+,1},\mathcal T^{(R)}_{+,1})
 &=({\cal N}_L\widehat{\mb Q}_+(\wh{\mb k}_1),
    {\cal N}_R\widehat{\mb Q}_+(\wh{\mb k}_3)),\nonumber\\
 (\mathcal T^{(L)}_{+,7},\mathcal T^{(R)}_{+,7})
 &=({\cal N}_L\mb Q_+(\wh{\mb k}_1),
    {\cal N}_R\ob{\mb Q}_+(\wh{\mb k}_3)),\nonumber\\
 (\mathcal T^{(L)}_{+,10},\mathcal T^{(R)}_{+,10})
 &=({\cal N}_L\ob{\mb Q}_+(\wh{\mb k}_1),
    {\cal N}_R\mb Q_+(\wh{\mb k}_3)),\nonumber\\
 (\mathcal T^{(L)}_{+,16},\mathcal T^{(R)}_{+,16})
 &=({\cal N}_L\check{\mb Q}_+(\wh{\mb k}_1),
    {\cal N}_R\check{\mb Q}_+(\wh{\mb k}_3)).
 \label{eq:four-hard-subgraphs}
\end{align}
Thus the subdiagrams in Eqs.~\eqref{eq:four-symmetric-traces} are now finite sums of the
You integral seeds, rather than unevaluated time integrals.  The type-1 and type-2
matrices $\mb I^D$ and
$\mb I^{\widehat D}$ are displayed explicitly in
Eqs.~\eqref{eq:seed-matrices} and \eqref{eq:hat-seed-matrices}.  The other two contain no
new integral,
\begin{align}
 \mb I^{\bar D;p_1p_2}_{ab}(\wh{\mb k})
 &=-\left[\mb I^{D;p_2p_1}_{ba}(-\wh{\mb k})\right]^T,\nonumber\\
 \mb I^{\check D;p_1p_2}_{ab}(\wh{\mb k})
 &=-\left[\mb I^{\widehat D;p_1^*p_2^*}_{-a,-b}
 (\wh{\mb k})\right]^* .
 \label{eq:conjugate-seeds-34}
\end{align}
After the index raising in Eq.~\eqref{eq:raised-hard-blocks}, the radial coefficients are
obtained without a further integral:
\begin{align}
 q_{X0}&=\frac12\left[(\mb Q^X_+)_{11}+(\mb Q^X_+)_{22}\right],\nonumber\\
 q_{X1}&=\frac{\eta_X}{2}\left[(\mb Q^X_+)_{11}-(\mb Q^X_+)_{22}\right],
 \qquad
 \eta_D=\eta_{\widehat D}=1,\qquad
 \eta_{\bar D}=\eta_{\check D}=-1.
 \label{eq:q-from-seed}
\end{align}
In Eq.~\eqref{eq:q-from-seed}, $q_{D i}=q_i$,
$q_{\bar D i}=\bar q_i$, $q_{\widehat D i}=\widehat q_i$, and
$q_{\check D i}=\check q_i$.

\newpage
To make the special-function dependence explicit, all entries needed in
Eqs.~\eqref{eq:Yukawa-dressed-seed}--\eqref{eq:q-from-seed} reduce to
\begin{align}
 {\mathfrak K}^{p_1p_2}(x,y)
 &=e^{\frac{\ii\pi}{2}(p_1-p_2)}
 \frac{\prod_{j=1}^{2}\left[\Gamma(1+p_j-\ii\tn)
 \Gamma(1+p_j+\ii\tn)\right]}
 {2^{2+p_1+p_2}\Gamma(p_1+x)\Gamma(p_2+y)},
 \label{eq:K-seed-main}\\
 {\mathfrak I}^{p_1p_2}(x,y,z)
 &=\ii\,2^{-2-p_1-p_2}e^{-\frac{\ii\pi}{2}(p_1+p_2)}
 \pi\operatorname{csch}(2\pi\tn)\Gamma(2+p_1+p_2)\nonumber\\
 &\quad\times\left[
 \frac{e^{-\pi\tn}\Gamma(1+p_1-\ii\tn)
 \Gamma(2+p_1+p_2-2\ii\tn)}{\Gamma(x)}\right.\nonumber\\[-1mm]
 &\hspace{12mm}\left.\times{}_4\widetilde F_3\!\left(
 \begin{matrix}2+p_1+p_2,\ 1+p_1-\ii\tn,\ y,
 2+p_1+p_2-2\ii\tn\\
 2+p_1-\ii\tn,\ 3+p_1+p_2-z,\ 1-2\ii\tn
 \end{matrix};-1\right)-(\tn\rightarrow-\tn)\right],
 \label{eq:I-seed-main}\\
 {\mathfrak J}^{p_1p_2}(x,y,z)
 &=\ii\,2^{-2-p_1-p_2}e^{\frac{\ii\pi}{2}(p_1+p_2)}
 \pi\operatorname{csch}(2\pi\tn)\Gamma(2+p_1+p_2)\nonumber\\
 &\quad\times\left[
 \frac{e^{\pi\tn}\Gamma(1+p_1-\ii\tn)
 \Gamma(2+p_1+p_2-2\ii\tn)}{\Gamma(x)}\right.\nonumber\\[-1mm]
 &\hspace{12mm}\left.\times{}_4\widetilde F_3\!\left(
 \begin{matrix}2+p_1+p_2,\ 1+p_1-\ii\tn,\ y,
 2+p_1+p_2-2\ii\tn\\
 2+p_1-\ii\tn,\ 3+p_1+p_2-z,\ 1-2\ii\tn
 \end{matrix};-1\right)-(\tn\rightarrow-\tn)\right].
 \label{eq:J-seed-main}
\end{align}
Here ${}_4\widetilde F_3={}_4F_3/\prod_{j=1}^{3}\Gamma(b_j)$ is the regularized
hypergeometric function.  The finite list of arguments entering the type-1 block is
\begin{align}
 A_1&=\tn^2\mathfrak I(1+\ii\tn,1-\ii\tn,\ii\tn),&
 A_2&=\mathfrak I(\ii\tn,-\ii\tn,1+\ii\tn),\nonumber\\
 B_1&=\mathfrak I^{p_2p_1}(\ii\tn,-\ii\tn,1+\ii\tn),&
 B_2&=\tn^2\mathfrak I^{p_2p_1}(1+\ii\tn,1-\ii\tn,\ii\tn),\nonumber\\
 U_1&=\tn^2\mathfrak K(2,2),&U_2&=\mathfrak K(1,1),\nonumber\\
 V_1&=e^{-\ii\pi(p_1-p_2)}\mathfrak K(1,1),&
 V_2&=\tn^2e^{-\ii\pi(p_1-p_2)}\mathfrak K(2,2),
 \label{eq:type1-dependence-main}
\end{align}
where the superscript $p_1p_2$ is understood on every $\mathfrak I$ and $\mathfrak K$
unless it is shown explicitly.  For the type-2 block the corresponding list is
\begin{align}
 S_1&=\tn\mathfrak I(1+\ii\tn,1-\ii\tn,1+\ii\tn),&
 S_2&=\tn\mathfrak I(\ii\tn,-\ii\tn,\ii\tn),\nonumber\\
 S_3&=\tn\mathfrak J^{p_2p_1}(\ii\tn,-\ii\tn,\ii\tn),&
 S_4&=\tn\mathfrak J^{p_2p_1}(1+\ii\tn,1-\ii\tn,1+\ii\tn),\nonumber\\
 W_1&=\tn\mathfrak K(2,1),&W_2&=\tn\mathfrak K(1,2).
 \label{eq:type2-dependence-main}
\end{align}
Equations~\eqref{eq:seed-matrices} and \eqref{eq:hat-seed-matrices} specify where these
entries occur in each SK matrix.  Hence Eqs.~\eqref{eq:K-seed-main}--
\eqref{eq:type2-dependence-main} expose the complete mass dependence; no hard time
integral is left unevaluated in the symmetric result.

Summing Eqs.~\eqref{eq:four-symmetric-traces}, it is useful to define
\begin{align}
 {\cal R}_0(\tn)
 &=\widehat q_0^{\,2}+\check q_0^{\,2}+2q_0\bar q_0,\nonumber\\
 {\cal R}_1(\tn)
 &=\widehat q_1^{\,2}+\check q_1^{\,2}-2q_1\bar q_1.
 \label{eq:R01-symmetric}
\end{align}
The remaining bubble functions can also be written entirely in Gamma functions,
\begin{align}
 \ka_+&=\ii\left[\frac{\Gamma(-2\ii\tn)}
 {\Gamma(-\ii\tn)}\right]^2,\nonumber\\
 f_0(-\ii\tn,-\ii\tn)
 &=\frac{\Gamma(-\frac32-2\ii\tn)
 \Gamma^2(\frac32+\ii\tn)}
 {\Gamma(3+2\ii\tn)\Gamma^2(-\ii\tn)}.
 \label{eq:kappa-f0-explicit}
\end{align}
The complete symmetric amplitude for one cyclic ordering is therefore
\begin{keyeqn}
\begin{align}
 C^{\rm Box}_{(S)}(\tn,\theta_1,\theta_3,\phi_3)
 &=A^{(S)}_0(\tn)+A^{(S)}_{13}(\tn)c_{13},
 \label{eq:Cbox-S}\\
 A^{(S)}_0(\tn)
 &=\frac{2^{1+12\ii\tn}\ka_+^2}{(4\pi)^{3/2}}
 f_0(-\ii\tn,-\ii\tn){\cal R}_0(\tn),\nonumber\\
 A^{(S)}_{13}(\tn)
 &=\frac{2^{1+12\ii\tn}\ka_+^2}{(4\pi)^{3/2}}
 f_0(-\ii\tn,-\ii\tn){\cal R}_1(\tn).
 \label{eq:A01-symmetric}
\end{align}
\end{keyeqn}
Using the gauge-boson-box normalization, the corresponding four-point function is
\begin{keyeqn}
\begin{align}
 \left.\la\de h^4\ra_{\rm Box}^{\prime}\right|_S
 &=-N_cg_t^4\frac{H^4k_s^3}
 {16k_1^3k_2^3k_3^3k_4^3}
 \left(\frac{k_s^2}{4k_{12}k_{34}}\right)^{2\ii\tn}
 \left[A^{(S)}_0+A^{(S)}_{13}c_{13}\right]+{\rm c.c.}
 +\order\!\left(\frac{k_s}{k_i}\right).
 \label{eq:final-symmetric-fourpoint}
\end{align}
\end{keyeqn}
This is the fermionic counterpart of the gauge-boson box result: a mass-dependent
amplitude multiplying a complex power of the collapsed momentum ratio.  In contrast to
the vector case, the soft spinor deltas leave only the monopole and the term linear in
$c_{13}$ for a fixed cyclic ordering.  The phase $2^{12\ii\tn}$ combines the normalization
of the two soft propagators, the choice of hard scales $k_{12}/2$ and $k_{34}/2$, and the
factor $4$ in the clock argument; it is convention dependent, whereas the full expression
including its complex conjugate is not.

\section{Antisymmetric soft-line contribution}
\label{sec:antisym}

We next retain the antisymmetric part that was left for future work in
Ref.~\cite{you_rhn}.  At zero chemical potential, Eq.~(2.39) and Appendix~C.3 of that
reference give
\begin{align}
 \rho_f(\tn)&\equiv
 \left[\frac{\Gamma(-2\ii f\tn)}{\Gamma(-\ii f\tn)}\right]^2,
 \qquad \ka_f=\ii f\rho_f,\\
 \left[D_{ab}(k;\tau,\tau')\right]_{{\rm NLoc},f}
 &=-\rho_f(4k^2\tau\tau')^{\ii f\tn}
 \Xi_{ab}(\tau,\tau')\,\wh{\mb k}\cdot\mb\si ,
 \label{eq:type1-exact}\\
 \Xi_{ab}(\tau,\tau')
 &=\frac{a+b}{2}\,\mathop{\rm sgn}(\tau-\tau')
 +\frac{b-a}{2},\qquad a,b=\pm1 .
 \label{eq:Xi}
\end{align}
Thus $\Xi_{-+}=1$, $\Xi_{+-}=-1$,
$\Xi_{++}=\mathop{\rm sgn}(\tau-\tau')$, and
$\Xi_{--}=-\mathop{\rm sgn}(\tau-\tau')$.  For the reversed type-1 line it is convenient
to keep the same coefficient $-\rho_f$ and use
$\wh{\mb k}\cdot\bar{\mb\si}=-\wh{\mb k}\cdot\mb\si$.  The other two propagators obey
\begin{align}
 [\widehat D^{(A)}]_{\rm NLoc}=0,
 \qquad [\check D^{(A)}]_{\rm NLoc}=0
 \qquad(\widetilde\lambda=0).
 \label{eq:anti-type23-zero}
\end{align}
Consequently the antisymmetric contribution of rows $1,7,10,16$ is identically zero.
The remaining twelve rows contain one or two type-1 soft lines and must be retained.

\subsection{Loop momentum integral and tensor bubble subgraphs}

We follow the gauge-boson box calculation and first do the soft loop momentum.  Put
$\mb K=\mb q+\mb k_s$ and $x=-\ii f\tn$.  Before this integral, the universal
nonanalytic factor is
\begin{align}
 {\cal LOOP}^{\rm soft}_{f,\ell}
 &=\left[\prod_{j=1}^{4}(-\tau_j)^{\ii f\tn}\right]
 2^{4\ii f\tn}\int\frac{\di^3\mb q}{(2\pi)^3}
 q^{-2x}K^{-2x}\,{\cal R}_{23}\otimes{\cal R}_{41},
 \label{eq:anti-loop-before-q}
\end{align}
where ${\cal R}=\si^0$ for a type-2/type-3 line and is
$\wh{\mb q}\cdot\mb\si$ or $\wh{\mb K}\cdot\mb\si$ for a type-1 line.  A reversed
type-1 line carries $\bar{\mb\si}=-\mb\si$; this sign will be denoted by $\zeta=-1$.
Only scalar, vector, and rank-two tensor moments occur.  Using the same loop seeds as the
gauge-boson calculation, define
\begin{align}
 f_c(a,b)&=\frac12f_0\left(a+\frac12,b-1\right)
 -\frac12f_0\left(a-\frac12,b\right)
 -\frac12f_0\left(a+\frac12,b\right),\\
 f_{cc}(a,b)&=\frac14f_0(a+1,b-2)+\frac14f_0(a-1,b)
 +\frac14f_0(a+1,b)\nonumber\\
 &\quad-\frac12f_0(a,b-1)-\frac12f_0(a+1,b-1)
 +\frac12f_0(a,b).
 \label{eq:fc-fcc}
\end{align}
and
\begin{align}
 {\cal V}_{q}(x)&=f_c(x,x),\qquad {\cal V}_{K}(x)=-f_c(x,x),\\
 {\cal J}_{0}(x)&=f_0\left(x-\frac12,x+\frac12\right)
 -\frac12f_0\left(x+\frac12,x+\frac12\right),\\
 {\cal J}_{L}(x)&=f_{cc}\left(x-\frac12,x+\frac12\right)
 +f_c\left(x,x+\frac12\right),\qquad
 {\cal J}_{T}(x)=\frac{{\cal J}_{0}(x)-{\cal J}_{L}(x)}{2}.
 \label{eq:soft-seeds-all}
\end{align}
After extracting $k_s^{3+4\ii f\tn}/(4\pi)^{3/2}$, the required moments are
\begin{align}
 \int\wh q_i&={\cal V}_{q}\wh k_{s\,i},\qquad
 \int\wh K_i={\cal V}_{K}\wh k_{s\,i},\nonumber\\
 \int\wh K_i\wh q_j
 ={\cal J}_{T}(\de_{ij}-\wh k_{s\,i}\wh k_{s\,j})
 +{\cal J}_{L}\wh k_{s\,i}\wh k_{s\,j}.
 \label{eq:soft-tensors}
\end{align}
The integral signs in Eq.~\eqref{eq:soft-tensors} include the common weights
$q^{-2x}K^{-2x}$.  The tensor trace gives the useful check
$2{\cal J}_T+{\cal J}_L={\cal J}_0$.  Restoring the factors stripped above,
\begin{align}
 {\cal LOOP}^{\rm soft}_{f,\ell}
 =\left[\prod_{j=1}^{4}(-\tau_j)^{\ii f\tn}\right]
 \frac{2^{4\ii f\tn}k_s^{3+4\ii f\tn}}{(4\pi)^{3/2}}\,
 {\mathfrak B}_{f,\ell}\times(\text{SK/time kernels}),
 \label{eq:anti-loop-after-q}
\end{align}
where $\mathfrak B_{f,\ell}$ is the vector or tensor in
Eq.~\eqref{eq:soft-tensors}.  Thus the soft loop momentum is completely integrated
before treating its time-ordering kernel.

\subsection{The sixteen rows in the antisymmetric sector}

The following table supplies the information needed for every row.  The pair in the third
column is $(X_L,X_R)$ for the hard edges $(12)$ and $(34)$.  A symbol
$D_A(K)$ or $\bar D_A(q)$ denotes a type-1 antisymmetric soft line and fixes both its
momentum and its sign $\zeta$, with $\zeta_D=+1$ and $\zeta_{\bar D}=-1$.
\begin{center}
\small
\begin{tabular}{cclll}
\toprule
$\ell$ & word & $(X_L,X_R)$ & edge $(23)$ & edge $(41)$\\
\midrule
1  & $NNNN$                     & $(\widehat D,\widehat D)$ & $\widehat D_S$ & $\widehat D_S$\\
2  & $NNN\bar N$                & $(\widehat D,D)$          & $\widehat D_S$ & $\bar D_A(q)$\\
3  & $NN\bar NN$                & $(\widehat D,\bar D)$     & $D_A(K)$       & $\widehat D_S$\\
4  & $NN\bar N\bar N$          & $(\widehat D,\check D)$   & $D_A(K)$       & $\bar D_A(q)$\\
5  & $N\bar NNN$                & $(D,\widehat D)$          & $\bar D_A(K)$  & $\widehat D_S$\\
6  & $N\bar NN\bar N$          & $(D,D)$                    & $\bar D_A(K)$  & $\bar D_A(q)$\\
7  & $N\bar N\bar NN$          & $(D,\bar D)$               & $\check D_S$   & $\widehat D_S$\\
8  & $N\bar N\bar N\bar N$    & $(D,\check D)$             & $\check D_S$   & $\bar D_A(q)$\\
9  & $\bar NNNN$                & $(\bar D,\widehat D)$     & $\widehat D_S$ & $D_A(q)$\\
10 & $\bar NNN\bar N$          & $(\bar D,D)$               & $\widehat D_S$ & $\check D_S$\\
11 & $\bar NN\bar NN$          & $(\bar D,\bar D)$         & $D_A(K)$       & $D_A(q)$\\
12 & $\bar NN\bar N\bar N$    & $(\bar D,\check D)$       & $D_A(K)$       & $\check D_S$\\
13 & $\bar N\bar NNN$          & $(\check D,\widehat D)$   & $\bar D_A(K)$  & $D_A(q)$\\
14 & $\bar N\bar NN\bar N$    & $(\check D,D)$             & $\bar D_A(K)$  & $\check D_S$\\
15 & $\bar N\bar N\bar NN$    & $(\check D,\bar D)$       & $\check D_S$   & $D_A(q)$\\
16 & $\bar N\bar N\bar N\bar N$ & $(\check D,\check D)$ & $\check D_S$   & $\check D_S$\\
\bottomrule
\end{tabular}
\end{center}
Rows $1,7,10,16$ contain no antisymmetric line and hence contribute zero to the present
sector.  The remaining rows fall into
\begin{align}
 {\cal G}_{AN}&=\{2,8,9,15\},&
 {\cal G}_{NA}&=\{3,5,12,14\},&
 {\cal G}_{NN}&=\{4,6,11,13\},
 \label{eq:anti-classes}
\end{align}
where $AN$ has a type-1 line on $(41)$, $NA$ has one on $(23)$, and $NN$ has one on
each soft edge.  This accounts for all $4+4+4+4=16$ rows.

\subsection{Time kernels and energy-deformed You seeds}

The sign kernel is separated by the exact spectral identity
\begin{align}
 \mathop{\rm sgn}(\tau-\tau')
 =\frac{1}{\ii\pi}\,{\rm PV}\int_{-\infty}^{+\infty}
 \frac{\di\omega}{\omega}\,e^{\ii\omega(\tau-\tau')}.
 \label{eq:sign-spectral}
\end{align}
For each fixed pair of SK labels introduce
\begin{align}
 C^{(0)}_{uv}=\frac{v-u}{2},\qquad
 C^{(1)}_{uv}=\frac{u+v}{2},\qquad
 \Xi_{uv}=C^{(0)}_{uv}+C^{(1)}_{uv}\mathop{\rm sgn}(\tau-\tau').
 \label{eq:C01}
\end{align}
Thus an opposite-SK type-1 line selects the algebraic term $C^{(0)}$, while an equal-SK
line selects one Hilbert transform through $C^{(1)}$.  This observation avoids any
double counting of SK sums.

The hard blocks must therefore be kept branch resolved.  Define the energy-deformed
Yukawa seed, without an SK sum, by
\begin{align}
 \mb Q^X_{f;ab}(\wh{\mb k};u,v)
 \equiv\int_0^\infty\di z_1\di z_2\,
 z_1^{p_f}z_2^{p_f}
 (1+\ii az_1)(1+\ii bz_2)
 e^{-\ii az_1-\ii bz_2+\ii uz_1+\ii vz_2}
 \mb D^X_{ab}(\wh{\mb k};z_1,z_2),
 \label{eq:deformed-seed}
\end{align}
where $p_f=-1+\ii f\tn$.  It is important that Eq.~\eqref{eq:deformed-seed} contains
no factor $ab$; the four branch labels are summed only after the cross-edge kernels have
been inserted.  The dependence on the special functions of Ref.~\cite{you_rhn} is made
explicit by the convergent, $\ii\epsilon$-regulated expansion
\begin{align}
 \mb Q^X_{f;ab}(u,v)
 =\sum_{r,s=0}^{\infty}\frac{(\ii u)^r(\ii v)^s}{r!s!}\Big[
 &\mb I^{X;p_f+r,p_f+s}_{ab}
 +\ii a\mb I^{X;p_f+r+1,p_f+s}_{ab}\nonumber\\[-1mm]
 &+\ii b\mb I^{X;p_f+r,p_f+s+1}_{ab}
 -ab\mb I^{X;p_f+r+1,p_f+s+1}_{ab}\Big].
 \label{eq:deformed-seed-series}
\end{align}
Every matrix $\mb I^{X;p_1p_2}_{ab}$ on the right-hand side is one of the You seeds in
Appendix~\ref{app:seeds}; hence every term is an explicit combination of Gamma functions
and ${}_4\widetilde F_3(-1)$.  Equation~\eqref{eq:deformed-seed-series}, analytically
continued from the regulated half-plane, is also a useful numerical representation of the
Hilbert transforms below.  At $u=v=0$ it reduces to the finite four-seed combination in
Eq.~\eqref{eq:Yukawa-dressed-seed} before the SK sum.

Let $\epsilon_{23},\epsilon_{41}\in\{0,1\}$ indicate whether the algebraic or sign term
in Eq.~\eqref{eq:C01} is selected.  Define
\begin{align}
 {\mathbb P}_{\epsilon_{23}\epsilon_{41}}F
 \equiv\prod_{e\in\{23,41\}:\epsilon_e=1}
 \left[\frac1{\ii\pi}{\rm PV}\!\int_{-\infty}^{+\infty}
 \frac{\di\Omega_e}{\Omega_e}\right]
 F\!\left(
 \frac{2\epsilon_{41}\Omega_{41}}{k_{12}},
 -\frac{2\epsilon_{23}\Omega_{23}}{k_{12}};
 \frac{2\epsilon_{23}\Omega_{23}}{k_{34}},
 -\frac{2\epsilon_{41}\Omega_{41}}{k_{34}}\right),
 \label{eq:PV-operator}
\end{align}
where an absent $\Omega_e$ is set to zero.  The four arguments of $F$ are respectively
$(u_L,v_L;u_R,v_R)$.  They follow directly from
$e^{\ii\Omega_{23}(\tau_2-\tau_3)}$ and
$e^{\ii\Omega_{41}(\tau_4-\tau_1)}$.

For row $\ell$, decompose the two branch-resolved hard blocks as
\begin{align}
 \mb Q^{X_L}_{+;ab}(\wh{\mb k}_1;u_L,v_L)
 &=L_{0;ab}^{\ell}\si^0+L_{1;ab}^{\ell}\wh{\mb k}_1\cdot\mb\si,\nonumber\\
 \mb Q^{X_R}_{+;cd}(\wh{\mb k}_3;u_R,v_R)
 &=R_{0;cd}^{\ell}\si^0+R_{1;cd}^{\ell}\wh{\mb k}_3\cdot\mb\si.
 \label{eq:branch-radial}
\end{align}
As in Eq.~\eqref{eq:q-from-seed}, the sign appropriate to a dotted block has already been
absorbed into $L_1$ or $R_1$.  For $m,n=0,1$, define the exact hard functional
\begin{align}
 {\cal H}^{mn}_\ell
 \equiv\sum_{a,b,c,d=\pm}abcd
 \sum_{\epsilon_{23},\epsilon_{41}}
 C^{(\epsilon_{23})}_{bc}C^{(\epsilon_{41})}_{da},
 {\mathbb P}_{\epsilon_{23}\epsilon_{41}}
 \left[L^{\ell}_{m;ab}R^{\ell}_{n;cd}\right].
 \label{eq:Hmn}
\end{align}
For a type-2/type-3 soft edge, the corresponding $\epsilon$ sum in
Eq.~\eqref{eq:Hmn} is replaced by the single value $\epsilon=0$ with $C^{(0)}=1$.
For a type-1 edge both values are retained.  Thus ${\cal H}^{mn}_\ell$ contains a finite
piece and at most one PV integral for $AN,NA$, and at most two PV integrals for $NN$.
Equations~\eqref{eq:deformed-seed-series}, \eqref{eq:PV-operator}, and
\eqref{eq:Hmn} are the corrected branch-resolved form of the antisymmetric master formula.
To write it uniformly, set
\begin{align}
 \chi^{\ell}_{e}=\begin{cases}
 \Xi_e,&\text{if edge $e$ is type-1 in row $\ell$},\\
 1,&\text{if edge $e$ is type-2/type-3 in row $\ell$}.
 \end{cases}
 \label{eq:chi-edge}
\end{align}
Here $\Sigma_e$ is the raised soft spinor matrix: it is $\si^0$ on a type-2/type-3 edge,
and $\zeta_e\widehat{\mb p}_e\cdot\mb\si$ on a type-1 edge, with
$\mb p_{23}=\mb K$ and $\mb p_{41}=\mb q$.  The bracket
$[\ \cdots\ ]_{\rm spectral}$ means that the vector or tensor bubble moment in
Eq.~\eqref{eq:soft-tensors} and the Hilbert operators in Eq.~\eqref{eq:PV-operator}
are applied before the branch sum.  Then
\begin{align}
 {\cal W}^{(A)}_{+,\ell}
 =\sum_{a,b,c,d=\pm}abcd\,
 \big[\chi^{\ell}_{23;bc}\chi^{\ell}_{41;da}\,
 \tr_2(\mb Q^{X_L}_{+;ab}\Sigma_{23}
       \mb Q^{X_R}_{+;cd}\Sigma_{41})\big]_{\rm spectral}.
 \label{eq:anti-master}
\end{align}

\subsection{Spinor contractions for the twelve nonzero rows}

Write
\begin{align}
 c_1&=\wh{\mb k}_1\cdot\wh{\mb k}_s=\cos\theta_1,
 &c_3&=\wh{\mb k}_3\cdot\wh{\mb k}_s=\cos\theta_3,\nonumber\\
 p_{13}&=\wh{\mb k}_s\cdot
 (\wh{\mb k}_1\times\wh{\mb k}_3)
 =\sin\theta_1\sin\theta_3\sin\phi_3.
 \label{eq:anti-angles}
\end{align}
For $\ell\in{\cal G}_{AN}$, the type-1 line is $(41)$ and the Pauli matrix lies to the
right of the right hard block.  Its contraction is
\begin{align}
 {\cal W}^{AN}_{\ell}
 =2\zeta_{41}^{\ell}{\cal V}_q\left(
 {\cal H}^{01}_\ell c_3+{\cal H}^{10}_\ell c_1
 +\ii{\cal H}^{11}_\ell p_{13}\right),
 \qquad \ell=2,8,9,15.
 \label{eq:WAN}
\end{align}
For $\ell\in{\cal G}_{NA}$, the type-1 matrix is between the two hard blocks, giving the
opposite sign for the three-Pauli trace,
\begin{align}
 {\cal W}^{NA}_{\ell}
 =2\zeta_{23}^{\ell}{\cal V}_K\left(
 {\cal H}^{01}_\ell c_3+{\cal H}^{10}_\ell c_1
 -\ii{\cal H}^{11}_\ell p_{13}\right),
 \qquad \ell=3,5,12,14.
 \label{eq:WNA}
\end{align}
Finally, contracting the symmetric tensor in Eq.~\eqref{eq:soft-tensors} with
$\tr(\mb Q_L\si^i\mb Q_R\si^j)$ gives
\begin{align}
 {\cal W}^{NN}_{\ell}
 =2\zeta_{23}^{\ell}\zeta_{41}^{\ell}\Big\{
 {\cal J}_0{\cal H}^{00}_\ell
 +{\cal H}^{11}_\ell\big[-{\cal J}_Lc_{13}
 +2({\cal J}_L-{\cal J}_T)c_1c_3\big]\Big\},
 \quad \ell=4,6,11,13.
 \label{eq:NN-angle}
\end{align}
All bubble functions in Eqs.~\eqref{eq:WAN}--\eqref{eq:NN-angle} are evaluated at
$x=-\ii\tn$.  The signs read directly from the table:
\begin{align}
 &(\zeta_{41}^{2},\zeta_{41}^{8},\zeta_{41}^{9},\zeta_{41}^{15})
 =(-1,-1,+1,+1),\nonumber\\
 &(\zeta_{23}^{3},\zeta_{23}^{5},\zeta_{23}^{12},\zeta_{23}^{14})
 =(+1,-1,+1,-1),\nonumber\\
 &(\zeta_{23}\zeta_{41})_{4,6,11,13}=(-1,+1,+1,-1).
 \label{eq:zeta-list}
\end{align}
Equations~\eqref{eq:WAN}--\eqref{eq:zeta-list}, together with the hard-pair column of the
table, are the individual result for each of the twelve nonzero rows.  No spinor index or
loop momentum remains uncontracted.

For clarity, the twelve answers can be displayed without any class notation.  Define
\begin{align}
 F^{\ell}_{\pm}
 &= {\cal H}^{01}_{\ell}c_3+{\cal H}^{10}_{\ell}c_1
 \mathbin{\pm}\ii {\cal H}^{11}_{\ell}p_{13},\nonumber\\
 G^{\ell}
 &= {\cal J}_0{\cal H}^{00}_{\ell}
 +{\cal H}^{11}_{\ell}\big[-{\cal J}_Lc_{13}
 +2({\cal J}_L-{\cal J}_T)c_1c_3\big].
 \label{eq:FG-anti}
\end{align}
The superscript $\ell$ fixes the hard pair listed in the sixteen-row table, so that, for
example, \({\cal H}^{mn}_2={\cal H}^{mn}[\widehat D,D]\) and
\({\cal H}^{mn}_{13}={\cal H}^{mn}[\check D,\widehat D]\).  The nonzero contractions are
\begin{align}
 &{\cal W}^{AN}_{2}=-2{\cal V}_qF^2_+,\qquad
 {\cal W}^{AN}_{8}=-2{\cal V}_qF^8_+,\qquad
 {\cal W}^{AN}_{9}=+2{\cal V}_qF^9_+,\qquad
 {\cal W}^{AN}_{15}=+2{\cal V}_qF^{15}_+,\nonumber\\
 &{\cal W}^{NA}_{3}=+2{\cal V}_KF^3_-,\qquad
 {\cal W}^{NA}_{5}=-2{\cal V}_KF^5_-,\qquad
 {\cal W}^{NA}_{12}=+2{\cal V}_KF^{12}_-,\qquad
 {\cal W}^{NA}_{14}=-2{\cal V}_KF^{14}_-,\nonumber\\
 &{\cal W}^{NN}_{4}=-2G^4,\qquad
 {\cal W}^{NN}_{6}=+2G^6,\qquad
 {\cal W}^{NN}_{11}=+2G^{11},\qquad
 {\cal W}^{NN}_{13}=-2G^{13}.
 \label{eq:twelve-anti-expanded}
\end{align}
Together with \({\cal W}^{(A)}_{1,7,10,16}=0\), Eq.~\eqref{eq:twelve-anti-expanded}
is an explicit sixteen-row audit of the antisymmetric sector.

\subsection{Nonlocal signal and angular coefficients}

One symmetric type-2/type-3 line contributes $\ka_+=\ii\rho_+$ and one antisymmetric
type-1 line contributes $-\rho_+$.  Two type-1 lines contribute $\rho_+^2$.  It follows
that
\begin{align}
 C^{\rm Box}_{(A)}
 ={\cal P}_A(\tn)\left[
 -\ii\sum_{\ell\in{\cal G}_{AN}}{\cal W}^{AN}_\ell
 -\ii\sum_{\ell\in{\cal G}_{NA}}{\cal W}^{NA}_\ell
 +\sum_{\ell\in{\cal G}_{NN}}{\cal W}^{NN}_\ell\right],
 \qquad
 {\cal P}_A(\tn)=\frac{2^{12\ii\tn}\rho_+^2}{(4\pi)^{3/2}}.
 \label{eq:Cbox-A}
\end{align}
This already has a form directly suitable for numerical evaluation.  To display the
angle dependence as in the gauge-boson box, write
\begin{align}
 C^{\rm Box}_{(A)}
 =A^{(A)}_0+A^{(A)}_1c_1+A^{(A)}_3c_3
 +A^{(A)}_{13}c_{13}+A^{(A)}_Lc_1c_3
 +A^{(A)}_Pp_{13}.
 \label{eq:Cbox-A-angular}
\end{align}
The six mass- and ratio-dependent coefficients are
\begin{subequations}\label{eq:anti-coefficients}
\begin{align}
 A^{(A)}_1
 &=-2\ii{\cal P}_A\left[
 \sum_{\ell\in{\cal G}_{AN}}\zeta_{41}^{\ell}{\cal V}_q{\cal H}^{10}_\ell
 +\sum_{\ell\in{\cal G}_{NA}}\zeta_{23}^{\ell}{\cal V}_K{\cal H}^{10}_\ell\right],\\
 A^{(A)}_3
 &=-2\ii{\cal P}_A\left[
 \sum_{\ell\in{\cal G}_{AN}}\zeta_{41}^{\ell}{\cal V}_q{\cal H}^{01}_\ell
 +\sum_{\ell\in{\cal G}_{NA}}\zeta_{23}^{\ell}{\cal V}_K{\cal H}^{01}_\ell\right],\\
 A^{(A)}_0
 &=2{\cal P}_A\sum_{\ell\in{\cal G}_{NN}}
 \zeta_{23}^{\ell}\zeta_{41}^{\ell}{\cal J}_0{\cal H}^{00}_\ell,\\
 A^{(A)}_{13}
 &=-2{\cal P}_A\sum_{\ell\in{\cal G}_{NN}}
 \zeta_{23}^{\ell}\zeta_{41}^{\ell}{\cal J}_L{\cal H}^{11}_\ell,\\
 A^{(A)}_L
 &=4{\cal P}_A\sum_{\ell\in{\cal G}_{NN}}
 \zeta_{23}^{\ell}\zeta_{41}^{\ell}
 ({\cal J}_L-{\cal J}_T){\cal H}^{11}_\ell,\\
 A^{(A)}_P
 &=2{\cal P}_A\left[
 \sum_{\ell\in{\cal G}_{AN}}\zeta_{41}^{\ell}{\cal V}_q{\cal H}^{11}_\ell
 -\sum_{\ell\in{\cal G}_{NA}}\zeta_{23}^{\ell}{\cal V}_K{\cal H}^{11}_\ell\right].
\end{align}
\end{subequations}
Here a sum labelled by a class runs over the four rows in Eq.~\eqref{eq:anti-classes}.
The coefficients depend on $\tn$ through $\rho_+$, the Gamma-function bubble seeds, and
the Gamma--${}_4\widetilde F_3$ series in Eq.~\eqref{eq:deformed-seed-series}.  They also
depend on the hard-scale ratio $r=k_{12}/k_{34}$ through the frequency arguments in
Eq.~\eqref{eq:PV-operator}.  Thus their mass and kinematic dependence is explicit.

Charge conjugation pairs the rows as
$(2,15)$, $(8,9)$, $(3,14)$, $(5,12)$, $(4,13)$, and $(6,11)$.  Applying
Eq.~\eqref{eq:conjugate-seeds-34} inside Eq.~\eqref{eq:Hmn}, followed where necessary by
the harmless PV change of variables $\Omega_e\to-\Omega_e$, shows that their parity-odd
pieces cancel pairwise, so
\begin{align}
 A^{(A)}_P=0.
 \label{eq:parity-cancel}
\end{align}
The exact antisymmetric answer for one fixed cyclic ordering is therefore the five-term
polynomial in Eq.~\eqref{eq:Cbox-A-angular} with its last term omitted.  This is the most
direct analogue of the gauge-boson angular decomposition.

The complete answer for a fixed cyclic ordering is
\begin{keyeqn}
\begin{align}
 \la\de h^4\ra_{\rm Box}^{\prime}
 &=-N_cg_t^4\frac{H^4k_s^3}
 {16k_1^3k_2^3k_3^3k_4^3}
 \left(\frac{k_s^2}{4k_{12}k_{34}}\right)^{2\ii\tn}
 C^{\rm Box}(\tn,r,\theta_1,\theta_3,\phi_3)+{\rm c.c.},
 \label{eq:final-fourpoint}\\
 C^{\rm Box}
 &=C^{\rm Box}_{(S)}+C^{\rm Box}_{(A)}\nonumber\\
 &=A_0+A_1c_1+A_3c_3+A_{13}c_{13}+A_Lc_1c_3,
 \label{eq:Cbox}
\end{align}
\end{keyeqn}
where
\begin{align}
 A_0&=A^{(S)}_0+A^{(A)}_0,&
 A_{13}&=A^{(S)}_{13}+A^{(A)}_{13},\nonumber\\
 A_1&=A^{(A)}_1,&A_3&=A^{(A)}_3,&A_L&=A^{(A)}_L.
\end{align}
The complex power is identical to the symmetric and gauge-boson boxes.  The only new
objects are one- and two-dimensional Hilbert transforms of known fermion seeds.  A finite
Gamma-function expression for these transforms is not available in
Ref.~\cite{you_rhn}; that paper explicitly identifies the antisymmetric time ordering as
the obstruction to direct cutting.  Equations~\eqref{eq:deformed-seed-series} and
\eqref{eq:PV-operator} isolate this obstruction without leaving any four-time integral.

For later reference, the symmetric angular result is
\begin{align}
 C^{\rm Box}_{(S)}=A^{(S)}_0+A^{(S)}_{13}c_{13}.
 \label{eq:AA-angle}
\end{align}
Combining it with Eq.~\eqref{eq:parity-cancel}, the angular dependence of one cyclic box is
\begin{align}
 C^{\rm Box}_{\rm cyclic}
 =A_0+A_{13}c_{13}+A_1c_1+A_3c_3+A_{L}c_1c_3.
 \label{eq:cyclic-angle}
\end{align}
There is no $\cos2\phi_3$ term.

The observable correlator of identical Higgs fluctuations also includes crossed cyclic
orderings.  At leading order, exchanging the two legs in the left pair sends
$\wh{\mb k}_1\to-\wh{\mb k}_1$, and the analogous exchange on the right sends
$\wh{\mb k}_3\to-\wh{\mb k}_3$.  Averaging independently over these exchanges projects
Eq.~\eqref{eq:cyclic-angle} onto
\begin{align}
 C^{\rm Box}_{\rm Bose}=A_0(\tn,r)+\order(k_s/k_i).
 \label{eq:Bose-angle}
\end{align}
Hence the leading, fully Bose-symmetrized top-box clock signal is isotropic.  The angular
terms in Eq.~\eqref{eq:cyclic-angle} belong to an ordered box and must not be confused
with the final identical-scalar angular dependence.

The scale dependence in the gauge-boson convention is
\begin{align}
 k_s^3\left(\frac{k_s^2}{4k_{12}k_{34}}\right)^{2\ii\tn}
 +{\rm c.c.}
 \ \sim\ k_s^3\cos\left[
 2\tn\log\left(\frac{k_s^2}{4k_{12}k_{34}}\right)+\varphi\right].
\end{align}

\section{Independent checks}

We close with checks that do not reuse the factorized derivation.

First, substituting Eq.~\eqref{eq:takagi} directly into the Dirac action reproduces both the
mass term and the Yukawa vertex in Eq.~\eqref{eq:maj-L}.  The transformation is unitary,
so its Grassmann Jacobian is one.  The two Majorana determinants reproduce one Dirac
determinant, confirming the $N_c$, rather than $2N_c$, normalization.

Second, the lower-index type-2/type-3 helicity sum is the invariant tensor $\ep$.  Raising one
index gives the identity, proving the rotation independence and the SK equality used in
the cut.  This also shows directly that no loop-momentum angular integral was silently
dropped in Eq.~\eqref{eq:soft-bubble}.

Third, enumerating the two vertex choices as four-bit words gives exactly sixteen entries.
The edge rule in Eq.~\eqref{eq:edge-rule} gives four entries in each of the $AA$, $AN$,
$NA$, and $NN$ classes, so no chirality flow is duplicated or omitted.

Fourth, the conformal loop formula gives $k_s^{3+4\ii f\tn}$.  The two hard subgraphs give
\begin{align}
 \left(\frac{k_{12}}{2}\right)^{-2\ii f\tn}
 \left(\frac{k_{34}}{2}\right)^{-2\ii f\tn},
\end{align}
so
Eq.~\eqref{eq:final-fourpoint} has the stated nonanalytic ratio and overall momentum
dimension $k^{-9}$.

Fifth, Eqs.~\eqref{eq:Xi} and \eqref{eq:sign-spectral} were checked distributionally:
$d\,\mathop{\rm sgn}(t)/dt=2\de(t)$ and the principal-value transform gives the same
derivative.  This independently fixes the signs of all four SK kernels.

Sixth, multiplying the four pairs of $2\times2$ matrices in the
$\mb k\parallel\wh{\mb z}$ frame and then applying the SU(2) rotations reproduces all
four traces in Eq.~\eqref{eq:four-symmetric-traces}.  This independently leaves only
$1$ and $c_{13}$ in the symmetric sector.  For the antisymmetric sector, direct
Pauli-matrix multiplication gives Eqs.~\eqref{eq:WAN}, \eqref{eq:WNA}, and
\eqref{eq:NN-angle}, including the opposite signs of the two triple-product terms.
Pairing every word with its charge conjugate then eliminates that triple product.  The tensor identity
$2{\cal J}_T+{\cal J}_L={\cal J}_0$ follows by tracing
Eq.~\eqref{eq:soft-tensors}; it also follows independently from
$\wh{\mb q}\cdot\wh{\mb K}=(q^2+K^2-k_s^2)/(2qK)$.

Seventh, multiplying the two late-time soft factors directly gives
\begin{align}
 (4q^2\tau_2\tau_3)^{\ii f\tn}
 (4K^2\tau_4\tau_1)^{\ii f\tn}
 =2^{4\ii f\tn}q^{2\ii f\tn}K^{2\ii f\tn}
 \prod_{j=1}^{4}(-\tau_j)^{\ii f\tn}.
\end{align}
This verifies independently the separation into the bubble subgraph and the four extra
time powers in Eqs.~\eqref{eq:symmetric-loop-after-q} and
\eqref{eq:anti-loop-after-q}.

Finally,
\begin{align}
 |\ka_+(\tn)|
 =\left|\frac{\Gamma(-2\ii\tn)}{\Gamma(-\ii\tn)}\right|^2
 =\frac{1}{4\cosh(\pi\tn)},
\end{align}
so the two cut soft lines contain the expected large-mass suppression
$|\ka_+|^2\sim e^{-2\pi\tn}/4$.  Moreover $\ka_-=\ka_+^*$,
$\mathcal Q_-=\mathcal Q_+^*$, and
$f_0(\ii\tn,\ii\tn)=f_0(-\ii\tn,-\ii\tn)^*$, which proves that the $f=-$ branch
is the complex conjugate displayed in Eq.~\eqref{eq:final-fourpoint}.

As a final seed check, setting $\widetilde\lambda=0$ and
$\widetilde M_R=\widetilde\nu$ in Eqs.~(2.52), (2.54), and (2.58) of
Ref.~\cite{you_rhn} gives Eqs.~\eqref{eq:K-seed-main}--\eqref{eq:J-seed-main}
term by term.  The reversal and complex-conjugation identities in
Eq.~\eqref{eq:conjugate-seeds-34} then reconstruct the $\bar D$ and $\check D$ blocks,
so the four symmetric diagrams require no seed beyond the three functions
$\mathfrak I$, $\mathfrak J$, and $\mathfrak K$.  Expanding the energy deformation as in
Eq.~\eqref{eq:deformed-seed-series} proves the same statement term by term for all twelve
nonzero antisymmetric diagrams.

\clearpage
\appendix

\section{Closed form of the hard fermion seed}
\label{app:seeds}

For completeness, this appendix collects the type-1, type-2, and type-3 hard seeds at
zero chemical potential, using the integral formulae of Ref.~\cite{you_rhn}.  They make
the four surviving terms in Eq.~\eqref{eq:Cbox-S} explicit in Gamma and generalized
hypergeometric functions.  The antisymmetric contribution remains in the exact
principal-value form of Eq.~\eqref{eq:anti-master}.

For $\mb k\parallel\wh{\mb z}$, write
\begin{align}
 \mathcal I_{++}^{p_1p_2}
 &=\begin{pmatrix}
 A_1^{p_1p_2}-B_2^{p_1p_2}&0\\
 0&A_2^{p_1p_2}-B_1^{p_1p_2}
 \end{pmatrix},
 \\
 \mathcal I_{--}^{p_1p_2}
 &=\left(\mathcal I_{++}^{p_2^*p_1^*}\right)^\dagger,
 \\
 \mathcal I_{-+}^{p_1p_2}
 &=\begin{pmatrix}U_1^{p_1p_2}&0\\0&U_2^{p_1p_2}\end{pmatrix},
 \qquad
 \mathcal I_{+-}^{p_1p_2}
 =\begin{pmatrix}-V_2^{p_1p_2}&0\\0&-V_1^{p_1p_2}\end{pmatrix}.
 \label{eq:seed-matrices}
\end{align}
The entries are
\begin{align}
 A_1^{p_1p_2}&=\tn^2\,
 \mathfrak I^{p_1p_2}(1+\ii\tn,1-\ii\tn,\ii\tn),\\
 A_2^{p_1p_2}&=
 \mathfrak I^{p_1p_2}(\ii\tn,-\ii\tn,1+\ii\tn),\\
 B_1^{p_1p_2}&=
 \mathfrak I^{p_2p_1}(\ii\tn,-\ii\tn,1+\ii\tn),\\
 B_2^{p_1p_2}&=\tn^2\,
 \mathfrak I^{p_2p_1}(1+\ii\tn,1-\ii\tn,\ii\tn),\\
 U_1^{p_1p_2}&=\tn^2\mathfrak K^{p_1p_2}(2,2),&
 U_2^{p_1p_2}&=\mathfrak K^{p_1p_2}(1,1),\\
 V_1^{p_1p_2}&=e^{-\ii\pi(p_1-p_2)}\mathfrak K^{p_1p_2}(1,1),&
 V_2^{p_1p_2}&=\tn^2e^{-\ii\pi(p_1-p_2)}\mathfrak K^{p_1p_2}(2,2).
\end{align}
The factorized seed is
\begin{align}
 \mathfrak K^{p_1p_2}(x,y)
 =e^{\frac{\ii\pi}{2}(p_1-p_2)}
 \frac{\Gamma(1+p_1-\ii\tn)\Gamma(1+p_1+\ii\tn)
 \Gamma(1+p_2-\ii\tn)\Gamma(1+p_2+\ii\tn)}
 {2^{2+p_1+p_2}\Gamma(p_1+x)\Gamma(p_2+y)}.
 \label{eq:K-seed}
\end{align}
The nested seed is
\begin{align}
 \mathfrak I^{p_1p_2}(x,y,z)
 &=\ii\,2^{-2-p_1-p_2}e^{-\frac{\ii\pi}{2}(p_1+p_2)}
 \pi\,\operatorname{csch}(2\pi\tn)\Gamma(2+p_1+p_2)
 \nonumber\\
 &\quad\times\Bigg[
 \frac{e^{-\pi\tn}}{\Gamma(x)}
 \Gamma(1+p_1-\ii\tn)
 \Gamma(2+p_1+p_2-2\ii\tn)
 \nonumber\\
 &\qquad\times{}_4\widetilde F_3\!\left(
 \begin{matrix}
 2+p_1+p_2,\ 1+p_1-\ii\tn,\ y,
 2+p_1+p_2-2\ii\tn\\
 2+p_1-\ii\tn,\ 3+p_1+p_2-z,\ 1-2\ii\tn
 \end{matrix};-1\right)
 -(\tn\rightarrow-\tn)\Bigg].
 \label{eq:I-seed}
\end{align}
Here ${}_4\widetilde F_3$ denotes the regularized generalized hypergeometric function.
Equations \eqref{eq:seed-matrices}--\eqref{eq:I-seed}, inserted into the
$rs=N\bar N$ case of Eq.~\eqref{eq:Q-def} with $p_+=-1+\ii\tn$, determine
$Q_{+,0}$ and $Q_{+,1}$ without any remaining time integral.  The reversed
$rs=\bar NN$ block follows from Eq.~\eqref{eq:Dbar}.

For the type-2 hard propagator, the corresponding lower-index matrices are
\begin{align}
 \widehat{\mathcal I}_{++}^{p_1p_2}
 &=\begin{pmatrix}0&S_1^{p_1p_2}+S_2^{p_2p_1}\\
 -(S_1^{p_2p_1}+S_2^{p_1p_2})&0\end{pmatrix},\\
 \widehat{\mathcal I}_{--}^{p_1p_2}
 &=\begin{pmatrix}0&S_3^{p_1p_2}+S_4^{p_2p_1}\\
 -(S_3^{p_2p_1}+S_4^{p_1p_2})&0\end{pmatrix},\\
 \widehat{\mathcal I}_{+-}^{p_1p_2}
 &=\begin{pmatrix}0&W_1^{p_1p_2}\\-W_2^{p_1p_2}&0\end{pmatrix},
 &
 \widehat{\mathcal I}_{-+}^{p_1p_2}
 &=\begin{pmatrix}0&W_2^{p_2p_1}\\-W_1^{p_2p_1}&0\end{pmatrix}.
 \label{eq:hat-seed-matrices}
\end{align}
At zero chemical potential their entries reduce to
\begin{align}
 S_1^{p_1p_2}&=\tn\,\mathfrak I^{p_1p_2}
 (1+\ii\tn,1-\ii\tn,1+\ii\tn),\\
 S_2^{p_1p_2}&=\tn\,\mathfrak I^{p_1p_2}
 (\ii\tn,-\ii\tn,\ii\tn),\\
 S_3^{p_1p_2}&=\tn\,\mathfrak J^{p_2p_1}
 (\ii\tn,-\ii\tn,\ii\tn),\\
 S_4^{p_1p_2}&=\tn\,\mathfrak J^{p_2p_1}
 (1+\ii\tn,1-\ii\tn,1+\ii\tn),\\
 W_1^{p_1p_2}&=\tn\,\mathfrak K^{p_1p_2}(2,1),&
 W_2^{p_1p_2}&=\tn\,\mathfrak K^{p_1p_2}(1,2).
 \label{eq:hat-seed-entries}
\end{align}
The second nested function is
\begin{align}
 \mathfrak J^{p_1p_2}(x,y,z)
 &=\ii\,2^{-2-p_1-p_2}e^{\frac{\ii\pi}{2}(p_1+p_2)}
 \pi\,\operatorname{csch}(2\pi\tn)\Gamma(2+p_1+p_2)
 \nonumber\\
 &\quad\times\Bigg[
 \frac{e^{\pi\tn}}{\Gamma(x)}
 \Gamma(1+p_1-\ii\tn)
 \Gamma(2+p_1+p_2-2\ii\tn)
 \nonumber\\
 &\qquad\times{}_4\widetilde F_3\!\left(
 \begin{matrix}
 2+p_1+p_2,\ 1+p_1-\ii\tn,\ y,
 2+p_1+p_2-2\ii\tn\\
 2+p_1-\ii\tn,\ 3+p_1+p_2-z,\ 1-2\ii\tn
 \end{matrix};-1\right)
 -(\tn\rightarrow-\tn)\Bigg].
 \label{eq:J-seed}
\end{align}
The type-3 seed follows without another integration,
\begin{align}
 \check{\mathcal I}_{ab}^{p_1p_2}
 =-\left(\widehat{\mathcal I}_{-a,-b}^{p_1^*p_2^*}\right)^*.
 \label{eq:check-seed}
\end{align}
Raising the appropriate second index and using
Eqs.~\eqref{eq:hat-seed-matrices}--\eqref{eq:check-seed} in
the corresponding $rs=NN$ and $rs=\bar N\bar N$ cases of Eq.~\eqref{eq:Q-def}
supplies the type-2 and type-3 hard blocks in
Eq.~\eqref{eq:Cbox-S}.  Thus all four entries of ${\cal G}_{AA}$ are explicit.

\begin{thebibliography}{9}

\bibitem{lu_chc}
S.~Lu, Y.~Wang, and Z.-Z.~Xianyu,
``A Cosmological Higgs Collider,''
JHEP \textbf{02} (2020) 011,
\href{https://arxiv.org/abs/1907.07390}{arXiv:1907.07390}.

\bibitem{you_rhn}
J.~You, L.~Song, C.~Han, H.-J.~He, X.~Chen, and Z.-Z.~Xianyu,
``Cosmological Collider Signatures from Right-Handed Neutrino Loop,''
\href{https://arxiv.org/abs/2605.21419}{arXiv:2605.21419} (2026).

\bibitem{qin_factorization}
Z.~Qin and Z.-Z.~Xianyu,
``Inflation Correlators at the One-Loop Order: Nonanalyticity, Factorization,
Cutting Rule, and OPE,'' JHEP \textbf{09} (2023) 116,
\href{https://arxiv.org/abs/2304.13295}{arXiv:2304.13295}.

\end{thebibliography}

\end{document}
